\chapter{Radial distribution function\label{App:Radial_distribution_function}}

Th radial distribution function, RDF or $g\left(r\right)$, is an elementary tool used to obtain structural information. \\
In chapter (\ref{chap7}) we present a numerical evaluation of the radial distribution function for systems with different correlations. In this appendix we show the used procedure. \\
The radial distribution function, from the numerical point of view, is determined by calculating the distance between all particle pairs and binning them into a histogram. The histogram is then normalized with respect to an ideal gas, where particle histograms are completely uncorrelated. For an uncorrelated system, the histogram can be expressed as follows.

Let us consider a first sphere with radius $R$, like schematically represented in Fig. (\ref{fig:app_sphere}). A second sphere, with radius $\rho$ and centred at distance $r$ from the centre of the first, is added to the system. 
%To evaluate the radial distribution we have to know the probability of find two particles at a distance $d$.
%The idea is to evaluate the of the surface 
%. This probability is given by the ratio between the surface area inside of the secondary inside the primary sphere by the total volume of the primary sphere. Let us consider a second sphere with center inside of the first sphere. The surface of this second sphere that is inside of the primary sphere is given by:
So, for any radius of the second sphere, the part of the surface that is embedded in the first sphere is given by:

\begin{equation}
S=2\pi\rho^{2}\int_{\theta_{2}}^{\pi}d\theta\sin\theta=2\pi\rho^{2}\left(1+\cos\theta_{2}\right).
\end{equation}

\begin{figure}[h!]
\begin{centering}
\includegraphics[width=8cm]{appendices/sphere_scheme.png} 
\par\end{centering}
\caption{Schematic representation of the sphere and its dimensions.\label{fig:app_sphere}}
\end{figure}
From the analysis of the Fig. (\ref{fig:app_sphere}), we can define the following equalities:

\begin{eqnarray}
R\sin\theta_{1} & = & \rho\sin\theta_{2} \nonumber\\
R\cos\theta_{1} & = & r+\rho\cos\theta_{2}
\end{eqnarray}
\begin{eqnarray}
R^{2}& = & \rho^{2}+r^{2}+2r\rho\cos\theta_{2}=\cos\theta_{2}=\frac{R^{2}-\rho^{2}-r^{2}}{2r\rho}\\
1+\cos\theta_{2}&=&\frac{R^{2}-\left(r-\rho\right)^{2}}{2r\rho}.
\end{eqnarray}

The probability of find a particle at some distance $r$ from the center of the first sphere is:

\begin{equation}
P\left(r\right)=\frac{4\pi r^2}{\left(4/3\right)\pi R^{3}}=\frac{3r^2}{R^3}.
\end{equation}

Now, the problem is divided in two parts. The first case, where $\left(\rho\leq R\right)$ and the second case, where $\left(\rho\geq R\right)$.

Case 1 $\left(\rho\leq R\right)$:

For the situation where $R\geq r + \rho$, \textit{i.e.}, the case where the second is totally inside of the primary sphere, the surface of the second sphere is given by: 

\begin{equation}
S\left(\rho,r\right)=4\pi\rho^{2}\qquad if\ R\geq r+\rho\rightarrow R-\rho\geq r\geq 0 \nonumber
\end{equation}

If only part of the second sphere is inside of the first sphere, $R<r+\rho$, the surface of second sphere embedded is:

\begin{equation}
S\left(\rho,r\right)=\pi\frac{\rho}{r}\left[R^{2}-\left(r-\rho\right)^{2}\right]\qquad R<r+\rho\rightarrow R-\rho<r.
\end{equation}

So, the probability of have any two particles at distance $\rho$ is given by:

\begin{eqnarray}
P\left(\rho\right) & \propto & \int_{0}^{R}S\left(\rho,r\right)P\left(r\right)dr\\
& \propto & \frac{3}{R^{3}}\left[\int_{0}^{R-\rho}dr4\pi r^{2}\rho^{2}+\int_{R-\rho}^{R}dr\pi\rho r\left[R^{2}-\left(r-\rho\right)^{2}\right]\right]\\
& \propto & \int_{0}^{R}S\left(\rho,r\right)P\left(r\right)dr=\frac{3}{R^{3}}\left[\int_{0}^{R-\rho}dr4\pi r^{2}\rho^{2}+\right.\nonumber\\
 & &\left.\int_{R-\rho}^{R}dr\pi\rho r\left[R^{2}-R^{2}-\rho^{2}+2r\rho-r^{2}\right]\right]\\
 & \propto & \int_{0}^{R}S\left(\rho,r\right)P\left(r\right)dr=\frac{3}{R^{3}}\left[\int_{0}^{R-\rho}dr4\pi r^{2}\rho^{2}+\right.\nonumber\\
 & &\left.\int_{R-\rho}^{R}dr\pi\rho r\left[R^{2}-R^{2}-\rho^{2}+2r\rho-r^{2}\right]\right]\\
 & \propto & \frac{3}{R^{3}}\left[\frac{4}{3}\pi\rho^{2}\left(R-\rho\right)^{3}+\frac{\pi\rho\left(R^{2}-\rho^{2}\right)}{2}\left[R^{2}-\left(R-\rho\right)^{2}\right]+\right.\nonumber\\
& & \left. \frac{2\pi\rho^{2}}{3}\left[R^{3}-\left(R-\rho\right)^{3}\right] -\frac{\pi\rho}{4}\left[R^{4}-\left(R-\rho\right)^{4}\right]\right]\\
P\left(\rho\right)& \propto & \frac{\pi}{4R^{3}}\left\{ 16\rho^{2}\left(R-\rho\right)^{3}+6\rho\left(R^{2}-\rho^{2}\right)\left[R^{2}-\left(R-\rho\right)^{2}\right]+\right.\nonumber\\
& &\left. +8\rho^{2}\left[R^{3}-\left(R-\rho\right)^{3}\right]-3\rho\left[R^{4}-\left(R-\rho\right)^{4}\right]\right\}.
\end{eqnarray}
%\begin{equation}
%If\ \rho=R\rightarrow P\left(\rho=R\right)=\frac{\pi}{4R^{3}}\left(8R^{5}-3R^{5}\right)=\frac{5}{4}\pi R^{2}
%\end{equation}

Case 2 $\left(\rho\geq R\right)$:

For the situation where $R\geq r - \rho$: 


 $\left(\rho\geq R\right)\ if\ \rho-r\leq R\rightarrow S\left(\rho,r\right)=\pi\frac{\rho}{R}\left[R^{2}-\left(r-\rho\right)^{2}\right]$

\[
\rho-r\geq R\rightarrow S\left(\rho,r\right)=0
\]


\[
P\left(\rho\right)\propto\frac{3\pi\rho}{R^{3}}\int_{\rho-R}^{R}\left\{ \left(R^{2}-\rho^{2}\right)r+2\rho r^{2}-r^{3}\right\} 
\]


\begin{eqnarray}
& \propto & \frac{3\pi\rho\left(R^{2}-\rho^{2}\right)}{2R^{3}}\left[R^{2}-\left(\rho-R\right)^{2}\right]+\nonumber\\
& & +\frac{2\pi\rho^{2}}{R^{3}}\left[R^{3}-\left(\rho-R\right)^{3}\right]-\frac{3\pi\rho}{4R^{3}}\left[R^{4}-\left(\rho-R\right)^{4}\right]
\end{eqnarray}


\begin{eqnarray}
P\left(\rho\right) & \propto &\frac{\pi\rho}{4R^{3}}\left\{ -3R^{4}+3\left(\rho-R\right)^{4}+8\rho\left[R^{3}-\left(\rho-R\right)^{3}\right]+\right.\nonumber\\
& & \left.+6\left(R^{2}-\rho^{2}\right)\left[R^{2}-\left(\rho-R\right)^{2}\right]\right\}.
\end{eqnarray}

So, we can finally write:

\begin{equation}
P\left(\rho\right) = \begin{cases}
 & \frac{3}{4\pi R^3}\frac{\pi}{4R^{3}}\left\{ 16\rho^{2}\left(R-\rho\right)^{3}+6\rho\left(R^{2}-\rho^{2}\right)\left[R^{2}-\left(R-\rho\right)^{2}\right]+\right.\\
&\left. +8\rho^{2}\left[R^{3}-\left(R-\rho\right)^{3}\right]-3\rho\left[R^{4}-\left(R-\rho\right)^{4}\right]\right\}  \quad \rho \le R\\
&\\ 
  &\frac{3}{4\pi R^3} \frac{\pi\rho}{4R^{3}}\left\{ -3R^{4}+3\left(\rho-R\right)^{4}+8\rho\left[R^{3}-\left(\rho-R\right)^{3}\right]+\right.\nonumber\\
& \left.+6\left(R^{2}-\rho^{2}\right)\left[R^{2}-\left(\rho-R\right)^{2}\right]\right\} \quad \rho \ge R.
\end{cases}
\end{equation}

In Fig. (\ref{fig:app_uncorrelated}) we present the normalized histogram of a system with radius $R = 1$.

\begin{figure}[h!]
\begin{centering}
\includegraphics[width=7cm, angle=-90]{appendices/uncorrelated.eps} 
\par\end{centering}
\caption{Normalised histogram of the distance between all particle pairs for a sphere with $R = 1$.\label{fig:app_uncorrelated}}
\end{figure}



%If $\rho=R^{+}\rightarrow P\left(\rho=R^{+}\right)=\frac{\pi R}{4R^{3}}\left\{ -3R^{4}+8R^{4}\right\} =\frac{5}{4}\pi R^{2}$

%That coincids with the other result $\rho=R^{-}$